\section{September 23: More Prime Numbers and Intro to Modular Arithmetic}

\subsection{More Results of Prime Numbers}

\begin{corrbox}
   Every integer $n$ can be written as $n = \sigma p_1^{a1} \dots p_k^{ak}$ for some $\sigma = \pm 1$, $k \geq 0$, and distinct primes $p_1 \dots p_k \in \NN$, and positive $a_1, \dots, a_k \in \NN$. Such a representation is unique up to reordering the $p_1 < \dots < p_k$. 
\end{corrbox}

This corrolary is just a restatement of the Fundamental Theorem of Arithmetic.

\begin{corrbox}
   Every rational $q = \sigma p_1^{a1} \dots p_k^{ak}$ for some $\sigma = \pm 1$, $k \geq 0$, distinct primes $p_1, \dots, p_k \in \NN$, and nonzero $a_1, \dots a_k \in \ZZ$. Suhc a representation is unique up to reordering the $p_1 < \dots < p_k$. 
\end{corrbox}

\begin{proof}
   The existence folows immediately from the existence of the corollary above by writing $q = m/n$ where $m \in \ZZ$ and $n \in \NN^+$. To show uniqueness, we prove the special case: if $1 = r_1{c1} \dots r_l^{cl}$ where $l \geq 0$, $r_1 \dots r_n$ are distinct prime numbers in $\NN$ and $c_1 \dots c_l \in \ZZ$ then for all $i \leq l$, $c_i = 0$. Indeed, if we suppose by contradiction that there exists an $i \leq l$ with $c_i \neq 0$, in particular $l \geq 1$. By removing those $r_j^{cj}$ where $c_j = 0$, we may assume without loss of generality that $c_i \neq 0$ for all $i \leq l$. By rearranging the $r_j$, we may assume without loss of generality that $c_1, \dots c_l < 0$ and $c_{l + }$
\end{proof}

<BS>
\begin{propbox}
   Any non-zero rational number $q$ can be written as
   \[
   q = \sigma p_1^{a_1} \dots p_m^{a_m},
   \]
   where $\sigma = \pm 1$, the $p_i$ are distinct prime numbers, and
   $a_1, \dots, a_m$ are non-zero integers, possibly negative. Moreover,
   this expression is unique up to reordering the primes.
\end{propbox}

\begin{proof}
   Since $q$ is rational, there exists $a$, $b \in \ZZ$ with $b \neq 0$ such that
   \[
   q = \frac{a}{b}.
   \]
   By the Fundamental Theorem of Arithmetic, we can write
   \[
   a = \sigma_a r_1^{s_1} \dots r_n^{s_n},
   \qquad
   b = \sigma_b r_1^{t_1} \dots r_n^{t_n},
   \]
   where $\sigma_a, \sigma_b \in \{-1,1\}$, the $r_i$ are distinct prime
   numbers, and $s_i, t_i$ are non-negative integers. We allow some of the
   exponents to be zero so that the same list of primes can be used in both
   factorizations. Therefore,
   \[
   q = \frac{\sigma_a}{\sigma_b}
   r_1^{s_1-t_1}\dots r_n^{s_n-t_n}.
   \]
   Omitting the primes for which $s_i-t_i=0$, and calling the remaining
   primes $p_1,\dots,p_m$, we obtain
   \[
   q = \sigma p_1^{a_1}\dots p_m^{a_m},
   \]
   where $\sigma=\sigma_a/\sigma_b \in \{-1,1\}$ and each $a_i$ is a
   non-zero integer.

   Now suppose that we have two such expressions for $q$. By allowing zero
   exponents, it is enough to consider
   \[
   q = \sigma p_1^{a_1}\dots p_m^{a_m}
   = \sigma' p_1^{a'_1}\dots p_m^{a'_m},
   \]
   where the $p_i$ are distinct prime numbers. Since $\sigma$ and $\sigma'$
   determine the sign of $q$, we must have $\sigma=\sigma'$. Dividing through
   by the second expression gives
   \[
   1 = p_1^{c_1}\dots p_m^{c_m},
   \]
   where $c_i=a_i-a'_i$. We need to show that $c_i=0$ for all $i$.

   By reordering the primes, we may assume that $c_1,\dots,c_t$ are
   negative and that $c_{t+1},\dots,c_m$ are non-negative. Hence,
   \[
   p_1^{-c_1}\dots p_t^{-c_t}
   =
   p_{t+1}^{c_{t+1}}\dots p_m^{c_m}.
   \]
   Both sides are now products of primes with non-negative exponents. By
   unique factorization for integers, and since all the primes are distinct,
   every exponent must be zero. Thus $c_i=0$ for all $i$, so
   \[
   a_i=a'_i
   \]
   for all $i$. Therefore, the expression is unique up to reordering the
   primes.
\end{proof}

